\chapter{Conclusiones}

\section{Conclusiones principales}

El trabajo demuestra que es viable estructurar la generación de modelos de
decisión como un pipeline de agentes especializados. La clave no es usar un
modelo de lenguaje de forma monolítica, sino dividir el proceso en etapas con
artefactos claros, revisión humana y validación ejecutable.

Las contribuciones principales son:

\begin{itemize}
\item Un Router determinista que coordina investigación, formalización,
razonamiento, construcción y memoria.
\item Un flujo de artefactos persistente que permite reanudar ejecuciones y
auditar decisiones.
\item Un contrato de modelos autocontenidos compatible con simulación dinámica.
\item Un sistema de memoria que combina PostgreSQL, MinIO, Neo4j y Qdrant.
\item Un mecanismo de recuperación híbrida con búsqueda densa, BM25, grafo,
fusión RRF y corrección tipo CRAG.
\item Una separación clara entre creatividad del modelo, control de sistema y
validación humana.
\end{itemize}

\section{Lecciones aprendidas}

La primera lección es que la memoria debe tener criterio de entrada. Guardar todo
lo que produce un modelo aumenta ruido y puede degradar ejecuciones futuras. Por
eso la arquitectura almacena conocimiento después de revisión.

La segunda es que la recuperación no puede depender de un único método. En un
dominio científico, nombres exactos, símbolos, autores y relaciones importan
tanto como similitud semántica. De ahí la combinación de BM25, embeddings y
grafo.

La tercera es que el código generado necesita pruebas, no solo validación textual.
Un modelo puede producir una explicación convincente y, aun así, escribir una
implementación que no respete el contrato de simulación.

\section{Trabajo futuro}

Las líneas futuras más relevantes son:

\begin{itemize}
\item Ampliar el catálogo de paradigmas de decisión y añadir referencias
específicas por experimento.
\item Mejorar métricas cuantitativas de calidad: precisión de slugs,
formulaciones válidas, tasa de tests superados y utilidad de memoria recuperada.
\item Añadir comparaciones sistemáticas entre modelos generados en entornos
equivalentes.
\item Extender la memoria con consultas temporales más ricas sobre la evolución
de formulaciones y modelos.
\item Integrar mejor la explicación de decisiones del modelo en los informes de
la fase 2.
\end{itemize}

\section{Cierre}

La fase 1 aporta la capa generativa del laboratorio virtual. Su valor está en
convertir conocimiento científico en modelos ejecutables manteniendo control
sobre las decisiones intermedias. Esta combinación de agentes, memoria y pruebas
ofrece una base sólida para estudiar paradigmas de decisión dentro de un entorno
de simulación reproducible.
